%% sol-ch04.tex — Solutions for Chapter 4: Nondeterministic Finite Automata

\chapter{Nondeterministic Finite Automata --- Solutions}

\begin{enumerate}[{4.}1.]

%% 4.1
\item Using the local Figure~4.25, the connected deterministic equivalents are
as follows.  In each table, $\emptyset$ is the dead state and all missing
subsets from $\wp(S)$ are disconnected.

\begin{enumerate}[(a)]
\item Figure~4.25(a): start state $\{t\}$, final state $\{t\}$.
\[
\begin{array}{c|ccc}
 & \pmb{a} & \pmb{b} & \pmb{c}\\\hline
\{t\} & \{t\} & \emptyset & \{t\}\\
\emptyset & \emptyset & \emptyset & \emptyset
\end{array}
\]

\item Figure~4.25(b): start state $\{q_1\}$, final state $\{q_1\}$.
\[
\begin{array}{c|ccc}
 & \pmb{a} & \pmb{b} & \pmb{c}\\\hline
\{q_1\} & \{q_0\} & \{q_1\} & \emptyset\\
\{q_0\} & \emptyset & \{q_1\} & \{q_0\}\\
\emptyset & \emptyset & \emptyset & \emptyset
\end{array}
\]

\item Figure~4.25(c): start state $\{r_0,r_1\}$, final states $\{r_0\}$ and
$\{r_0,r_1\}$.
\[
\begin{array}{c|ccc}
 & \pmb{a} & \pmb{b} & \pmb{c}\\\hline
\{r_0\} & \{r_0,r_1\} & \emptyset & \{r_1\}\\
\{r_0,r_1\} & \{r_0,r_1\} & \{r_0\} & \{r_1\}\\
\{r_1\} & \emptyset & \{r_0\} & \emptyset\\
\emptyset & \emptyset & \emptyset & \emptyset
\end{array}
\]

\item Figure~4.25(d): start state $\{s_0\}$, final state $\{s_2\}$.
\[
\begin{array}{c|ccc}
 & \pmb{a} & \pmb{b} & \pmb{c}\\\hline
\{s_0\} & \{s_1\} & \emptyset & \emptyset\\
\{s_1\} & \{s_2\} & \{s_3\} & \emptyset\\
\{s_2\} & \emptyset & \emptyset & \emptyset\\
\{s_3\} & \emptyset & \{s_1\} & \emptyset\\
\emptyset & \emptyset & \emptyset & \emptyset
\end{array}
\]

\item Figure~4.25(e): start state $\{s_1,s_2\}$, final state $\{s_1,s_2\}$.
\[
\begin{array}{c|ccc}
 & \pmb{a} & \pmb{b} & \pmb{c}\\\hline
\{s_1,s_2\} & \{s_1,s_2\} & \{s_1\} & \emptyset\\
\{s_1\} & \{s_1,s_2\} & \emptyset & \emptyset\\
\emptyset & \emptyset & \emptyset & \emptyset
\end{array}
\]

\item Figure~4.25(f): after taking the $\lambda$-closures, the connected DFA has
start state $\{s_0\}$ and final state $\{s_0,s_1,s_3\}$.
\[
\begin{array}{c|ccc}
 & \pmb{a} & \pmb{b} & \pmb{c}\\\hline
\{s_0\} & \{s_0,s_1\} & \emptyset & \emptyset\\
\{s_0,s_1\} & \{s_0,s_1\} & \emptyset & \{s_2\}\\
\{s_2\} & \emptyset & \emptyset & \{s_0,s_1,s_3\}\\
\{s_0,s_1,s_3\} & \{s_0,s_1\} & \emptyset & \{s_2\}\\
\emptyset & \emptyset & \emptyset & \emptyset
\end{array}
\]

\item Figure~4.25(g): start state $\{s_0\}$, final state $\{s_0\}$.
\[
\begin{array}{c|ccc}
 & \pmb{a} & \pmb{b} & \pmb{c}\\\hline
\{s_0\} & \{s_1\} & \emptyset & \emptyset\\
\{s_1\} & \{s_0\} & \{s_3\} & \{s_2\}\\
\{s_2\} & \emptyset & \emptyset & \{s_1\}\\
\{s_3\} & \emptyset & \{s_1\} & \emptyset\\
\emptyset & \emptyset & \emptyset & \emptyset
\end{array}
\]
\end{enumerate}

%% 4.2
\item Consider Example~4.17, that is, the automaton $E$ of Figure~4.23.
Reading the figure gives
\[
E=\langle\{\pmb{a},\pmb{b}\},\{s_1,s_2\},\{s_1\},\delta_{\lambda},\{s_2\}\rangle
\]
with ordinary transitions
\[
\delta_{\lambda}(s_1,\pmb{a})=\{s_1,s_2\},\qquad
\delta_{\lambda}(s_1,\pmb{b})=\emptyset,
\]
\[
\delta_{\lambda}(s_2,\pmb{a})=\{s_1\},\qquad
\delta_{\lambda}(s_2,\pmb{b})=\{s_1\},
\]
and one $\lambda$-transition $\delta_{\lambda}(s_1,\lambda)=\{s_2\}$.

\begin{enumerate}[(a)]
  \item \emph{Remove $\lambda$-transitions (Definition~4.9).}  The
  $\lambda$-closures are:
  \[
    \Lambda(s_1)=\{s_1,s_2\},\qquad
    \Lambda(s_2)=\{s_2\}.
  \]
  Using the $\lambda$-closure of the original start state, the corresponding
  $\lambda$-free NDFA has start set $\{s_1,s_2\}$. Thus
  \[
  E^{\sigma}_{\lambda}=\langle\{\pmb{a},\pmb{b}\},\{s_1,s_2\},\{s_1,s_2\},
  \delta^{\sigma}_{\lambda},\{s_2\}\rangle
  \]
  with
  \[
  \delta^{\sigma}_{\lambda}(s_1,\pmb{a})=\{s_1,s_2\},\qquad
  \delta^{\sigma}_{\lambda}(s_1,\pmb{b})=\{s_1,s_2\},
  \]
  \[
  \delta^{\sigma}_{\lambda}(s_2,\pmb{a})=\{s_1,s_2\},\qquad
  \delta^{\sigma}_{\lambda}(s_2,\pmb{b})=\{s_1,s_2\}.
  \]

  \item \emph{Convert this NDFA into a DFA (Definition~4.5).}  The full
  deterministic machine has $2^2=4$ subset states, but only one is reachable
  from the start set $\{s_1,s_2\}$:
  \[
  Q_0=\{s_1,s_2\}.
  \]
  This state is final because it contains $s_2$. Its
  transitions are
  \[
  \begin{array}{c|cc}
   & \pmb{a} & \pmb{b}\\\hline
  Q_0 & Q_0 & Q_0
  \end{array}
  \]
  so the connected part of the DFA is a one-state accepting machine.
\end{enumerate}

%% 4.3
\item Consider the NDFA $A$ from Example~4.4:
$A=\langle\{\pmb{a},\pmb{b}\},\{r,s,t\},\{r,s\},\delta,\{t\}\rangle$
with
\[
\delta(r,\pmb{a})=\{s\},\quad \delta(r,\pmb{b})=\emptyset,\quad
\delta(s,\pmb{a})=\emptyset,\quad \delta(s,\pmb{b})=\{r,t\},
\]
\[
\delta(t,\pmb{a})=\emptyset,\qquad \delta(t,\pmb{b})=\{t\}.
\]

\begin{enumerate}[(a)]
  \item \emph{Circuit diagram (no \textless SOS\textgreater{} or
    \textless EOS\textgreater).}  Using one activity bit for each state,
    the next-state equations are
    \[
    r' = s\wedge \pmb{b},\qquad
    s' = r\wedge \pmb{a},\qquad
    t' = (s\wedge \pmb{b}) \vee (t\wedge \pmb{b}).
    \]
    Initially the active states are $r$ and $s$, and $t$ is inactive.

  \item \emph{Circuit diagram (with \textless SOS\textgreater{} and
    \textless EOS\textgreater).}  \<SOS\> resets the activity bits to
    $r=s=1,\ t=0$.  The accept line is $t\wedge\text{\<EOS\>}$.

  \item \emph{DFA via subset construction (connected portion only).}
  The start state of $A^d$ is $\{r,s\}$.  The reachable states are
  \[
  \{r,s\},\quad \{s\},\quad \{r,t\},\quad \{t\},\quad \emptyset.
  \]
  Their transition table is
  \begin{center}\begin{tabular}{c|cc}
    State & $\pmb{a}$ & $\pmb{b}$ \\\hline
    $\{r,s\}$ & $\{s\}$ & $\{r,t\}$ \\
    $\{s\}$ & $\emptyset$ & $\{r,t\}$ \\
    $\{r,t\}$ & $\{s\}$ & $\{t\}$ \\
    $\{t\}$ & $\emptyset$ & $\{t\}$ \\
    $\emptyset$ & $\emptyset$ & $\emptyset$ \\
  \end{tabular}\end{center}
  The final states are $\{r,t\}$ and $\{t\}$.

  \item \emph{DFA including disconnected portion.}
  The full state set is
  \[
  \wp(\{r,s,t\})=\{\emptyset,\{r\},\{s\},\{t\},\{r,s\},\{r,t\},\{s,t\},\{r,s,t\}\}.
  \]
  Its transition table is
  \begin{center}\begin{tabular}{c|cc}
    State & $\pmb{a}$ & $\pmb{b}$ \\\hline
    $\emptyset$ & $\emptyset$ & $\emptyset$ \\
    $\{r\}$ & $\{s\}$ & $\emptyset$ \\
    $\{s\}$ & $\emptyset$ & $\{r,t\}$ \\
    $\{t\}$ & $\emptyset$ & $\{t\}$ \\
    $\{r,s\}$ & $\{s\}$ & $\{r,t\}$ \\
    $\{r,t\}$ & $\{s\}$ & $\{t\}$ \\
    $\{s,t\}$ & $\emptyset$ & $\{r,t\}$ \\
    $\{r,s,t\}$ & $\{s\}$ & $\{r,t\}$
  \end{tabular}\end{center}
  Final states are exactly those subsets containing $t$.

\end{enumerate}

%% 4.4
\item Consider the NDFA from Example~4.2, which has multiple transitions on
$\pmb{0}$: from $s_0$ to both $s_1$ and $s_2$.

\begin{enumerate}[(a)]
  \item \emph{Circuit diagram (with \textless SOS\textgreater{} and
    \textless EOS\textgreater).}  Use one activity flip-flop for each
    state $s_0,\dots,s_5$.  The next-state equations are
    \[
    s'_0=\text{\<SOS\>},\qquad
    s'_1=(s_0\wedge \pmb{0})\vee(s_1\wedge \pmb{0})\vee(s_5\wedge \pmb{1}),
    \]
    \[
    s'_2=(s_0\wedge \pmb{0})\vee(s_0\wedge \pmb{1})\vee
         (s_4\wedge \pmb{0})\vee(s_4\wedge \pmb{1}),
    \]
    \[
    s'_3=(s_0\wedge \pmb{1})\vee(s_1\wedge \pmb{1})\vee(s_3\wedge \pmb{0}),
    \]
    \[
    s'_4=(s_2\wedge \pmb{0})\vee(s_2\wedge \pmb{1}),\qquad
    s'_5=(s_3\wedge \pmb{1})\vee(s_5\wedge \pmb{0}).
    \]
    Since the final states in Figure~4.3 are $s_0$, $s_1$, and $s_4$, the
    accept line is
    \[
      (s_0\vee s_1\vee s_4)\wedge \text{\<EOS\>}.
    \]

  \item \emph{DFA via subset construction.}  Starting from the set of
  start state $\{s_0\}$, the connected portion of the subset DFA has the
  following reachable states:
  \[
  Q_0=\{s_0\},\quad
  Q_1=\{s_1,s_2\},\quad
  Q_2=\{s_2,s_3\},\quad
  Q_3=\{s_1,s_4\},\quad
  Q_4=\{s_3,s_4\},
  \]
  \[
  Q_5=\{s_4,s_5\},\quad
  Q_6=\{s_2,s_5\}.
  \]
  Their transitions are
  \[
  \begin{array}{c|cc}
   & \pmb{0} & \pmb{1}\\\hline
  Q_0 & Q_1 & Q_2\\
  Q_1 & Q_3 & Q_4\\
  Q_2 & Q_4 & Q_5\\
  Q_3 & Q_1 & Q_2\\
  Q_4 & Q_2 & Q_6\\
  Q_5 & Q_6 & Q_3\\
  Q_6 & Q_5 & Q_3
  \end{array}
  \]
  The final states are exactly those subsets containing one of the original
  final states $s_0$, $s_1$, or $s_4$, namely
  \[
  Q_0,\ Q_1,\ Q_3,\ Q_4,\ Q_5.
  \]
\end{enumerate}

%% 4.5
\item Consider the NDFA from Example~4.3, which accepts the same language as
Example~4.2 but uses multiple start states instead of multiple transitions.

\begin{enumerate}[(a)]
  \item \emph{Circuit diagram (no \textless SOS\textgreater{} or
    \textless EOS\textgreater).}  Again use one activity flip-flop per
    state.  For the automaton of Figure~4.4,
    \[
    s'_1=(s_1\wedge \pmb{0})\vee(s_5\wedge \pmb{0}),\qquad
    s'_2=(s_4\wedge \pmb{0})\vee(s_4\wedge \pmb{1}),
    \]
    \[
    s'_3=(s_1\wedge \pmb{1})\vee(s_3\wedge \pmb{0}),\qquad
    s'_4=(s_2\wedge \pmb{0})\vee(s_2\wedge \pmb{1}),
    \]
    \[
    s'_5=s_3\wedge \pmb{1}.
    \]
    The initial configuration has $s_1=s_4=1$ and the other activity bits
    equal to $0$.  Since the final states are $s_1$ and $s_4$, the accept
    line is $s_1\vee s_4$.

  \item \emph{Circuit diagram (with \textless SOS\textgreater{} and
    \textless EOS\textgreater).}  Keep the equations for $s'_2$, $s'_3$, and
    $s'_5$ from part~(a), keep $s'_0=0$, and modify only the two start-state
    inputs by OR-ing in the start pulse:
    \[
      s'_1=\text{\<SOS\>}\vee(s_1\wedge \pmb{0})\vee(s_5\wedge \pmb{0}),
      \qquad
      s'_4=\text{\<SOS\>}\vee(s_2\wedge \pmb{0})\vee(s_2\wedge \pmb{1}),
    \]
    with the equations for $s'_2$, $s'_3$, and $s'_5$ unchanged.
    The accept output is
    \[
      (s_1\vee s_4)\wedge \text{\<EOS\>}.
    \]

  \item \emph{DFA (connected portion only).}  The start state of $A^d$
  is the set of all start states, namely $\{s_1,s_4\}$.  The reachable
  subsets are
  \[
  P_0=\{s_1,s_4\},\quad
  P_1=\{s_1,s_2\},\quad
  P_2=\{s_2,s_3\},\quad
  P_3=\{s_3,s_4\},\quad
  P_4=\{s_4,s_5\},\quad
  P_5=\{s_2,s_5\},
  \]
  \[
  P_6=\{s_2\},\qquad
  P_7=\{s_4\}.
  \]
  Their transitions are
  \[
  \begin{array}{c|cc}
   & \pmb{0} & \pmb{1}\\\hline
  P_0 & P_1 & P_2\\
  P_1 & P_0 & P_3\\
  P_2 & P_3 & P_4\\
  P_3 & P_2 & P_5\\
  P_4 & P_1 & P_6\\
  P_5 & P_0 & P_7\\
  P_6 & P_7 & P_7\\
  P_7 & P_6 & P_6
  \end{array}
  \]
  The final states are
  \[
  P_0,\ P_1,\ P_3,\ P_4,\ P_7,
  \]
  since these are exactly the reachable subsets containing $s_1$ or $s_4$.

  \item \emph{Is this DFA isomorphic to any of those in Exercise~4.4?}
  No.  The DFA from Exercise~4.4(b) has seven reachable states
  \[
  Q_0,\ldots,Q_6,
  \]
  while the DFA from part~(c) above has eight reachable states
  \[
  P_0,\ldots,P_7.
  \]
  Since isomorphic automata must have the same number of states, these two
  subset-construction DFAs are not isomorphic.  They do accept the same
  language, but that only implies that their \emph{minimal} DFAs are
  isomorphic, not these unreduced machines.
\end{enumerate}

%% 4.6
\item Consider the $\lambda$-NDFA from Example~4.14, with states
$\{s_0,s_1,s_2,s_3\}$ and $\lambda$-transitions $s_0\to s_1$ and $s_1\to s_2$.

\begin{enumerate}[(a)]
  \item \emph{Circuit diagram for the original $\lambda$-NDFA (no SOS/EOS).}
  Figure~4.19 uses four activity bits, one for each state.  Ignoring the
  $\lambda$-moves for a moment, the ordinary transitions give
  \[
    s'_0=0,\qquad
    s'_1=(s_0\wedge \pmb{a})\vee(s_3\wedge \pmb{b}),
  \]
  \[
    s'_2=(s_1\wedge \pmb{c})\vee(s_2\wedge \pmb{d}),\qquad
    s'_3=s_1\wedge \pmb{b}.
  \]
  Because a $\lambda$-move consumes no input symbol, whenever $s_0$ is active
  the next-state logic must also make $s_1$ active, and whenever $s_1$ is
  active it must also make $s_2$ active.  This is achieved by OR-ing the source
  activity line directly into the destination input.  Thus the circuit can be
  specified by
  \[
    D_0=0,\qquad
    D_1=s_0\vee(s_0\wedge \pmb{a})\vee(s_3\wedge \pmb{b}),
  \]
  \[
    D_2=s_1\vee(s_1\wedge \pmb{c})\vee(s_2\wedge \pmb{d}),\qquad
    D_3=s_1\wedge \pmb{b}.
  \]
  Initially only $s_0$ is active; the accept output is simply $s_2$.

  \item \emph{Circuit diagram for the NDFA from Definition~4.9 (part~(b) of
  Example~4.14, no SOS/EOS).}  The corresponding $\lambda$-free automaton is
  the one shown in Figure~4.20.  Use one activity bit for each of
  $q_0,s_0,s_1,s_2,s_3$; then
  \[
    q'_0=0,\qquad s'_0=0,
  \]
  \[
    s'_1=(s_0\wedge \pmb{a})\vee(s_3\wedge \pmb{b}),
  \]
  \[
    s'_2=(s_0\wedge \pmb{a})\vee(s_0\wedge \pmb{c})\vee(s_0\wedge \pmb{d})
       \vee(s_1\wedge \pmb{c})\vee(s_1\wedge \pmb{d})
       \vee(s_3\wedge \pmb{b})\vee(s_2\wedge \pmb{d}),
  \]
  \[
    s'_3=(s_0\wedge \pmb{b})\vee(s_1\wedge \pmb{b}).
  \]
  The initial activity bits are $q_0=s_0=1$ and $s_1=s_2=s_3=0$, and the
  accept output is $q_0\vee s_2$.
\end{enumerate}

%% 4.7
\item Consider the second NDFA of Example~4.16, which accepts the language
$L^+=\{x_1 x_2\cdots x_k\mid k\ge 1,  x_i\in L\}$ where $L$ is the language
of strings over $\{\pmb{a},\pmb{b},\pmb{c}\}$ containing exactly one $\pmb{b}$
immediately followed by $\pmb{c}$.

\begin{enumerate}[(a)]
  \item \emph{Circuit diagram for $A^{\sigma}_\lambda$ (with EOS but no SOS).}
  For the $\lambda$-NDFA of Figure~4.22, use one activity bit for each of
  $s_0,s_1,s_2$.  The ordinary transitions give
  \[
    s'_0=s_0\wedge(\pmb{a}\vee\pmb{c}),\qquad
    s'_1=s_0\wedge \pmb{b},
  \]
  \[
    s'_2=(s_1\wedge \pmb{c})\vee(s_2\wedge(\pmb{a}\vee\pmb{c})).
  \]
  Because Figure~4.22 also has the $\lambda$-transition $s_2\to s_0$, OR the
  activity line for $s_2$ directly into the input for $s_0$.  Thus
  \[
    D_0=s_2\vee(s_0\wedge(\pmb{a}\vee\pmb{c})),\qquad
    D_1=s_0\wedge \pmb{b},
  \]
  \[
    D_2=(s_1\wedge \pmb{c})\vee(s_2\wedge(\pmb{a}\vee\pmb{c})).
  \]
  The accept line is $s_2\wedge\text{\<EOS\>}$.

  \item \emph{NDFA without $\lambda$-transitions (Definition~4.9).}
  For Figure~4.22 the $\lambda$-closures are
  \[
    \Lambda(s_0)=\{s_0\},\qquad
    \Lambda(s_1)=\{s_1\},\qquad
    \Lambda(s_2)=\{s_0,s_2\}.
  \]
  Since $\lambda\notin L(A_\lambda)$, Definition~4.9 adds $q_0$ as an extra
  disconnected start state but does not make it final.  Thus
  \[
    A^{\sigma}_\lambda=
    \langle\{\pmb{a},\pmb{b},\pmb{c}\},\{q_0,s_0,s_1,s_2\},\{q_0,s_0\},
    \delta^{\sigma}_\lambda,\{s_2\}\rangle,
  \]
  where
  \[
    \delta^{\sigma}_\lambda(q_0,\pmb{a})=
    \delta^{\sigma}_\lambda(q_0,\pmb{b})=
    \delta^{\sigma}_\lambda(q_0,\pmb{c})=\emptyset,
  \]
  \[
    \delta^{\sigma}_\lambda(s_0,\pmb{a})=\{s_0\},\qquad
    \delta^{\sigma}_\lambda(s_0,\pmb{b})=\{s_1\},\qquad
    \delta^{\sigma}_\lambda(s_0,\pmb{c})=\{s_0\},
  \]
  \[
    \delta^{\sigma}_\lambda(s_1,\pmb{a})=\emptyset,\qquad
    \delta^{\sigma}_\lambda(s_1,\pmb{b})=\emptyset,\qquad
    \delta^{\sigma}_\lambda(s_1,\pmb{c})=\{s_0,s_2\},
  \]
  \[
    \delta^{\sigma}_\lambda(s_2,\pmb{a})=\{s_0,s_2\},\qquad
    \delta^{\sigma}_\lambda(s_2,\pmb{b})=\{s_1\},\qquad
    \delta^{\sigma}_\lambda(s_2,\pmb{c})=\{s_0,s_2\}.
  \]

  \item \emph{Circuit diagram for the $\lambda$-free NDFA (with EOS, no SOS).}
  Use one activity bit for each of $q_0,s_0,s_1,s_2$.  The next-state
  equations are
  \[
    q'_0=0,
  \]
  \[
    s'_0=(s_0\wedge \pmb{a})\vee(s_0\wedge \pmb{c})
        \vee(s_1\wedge \pmb{c})\vee(s_2\wedge \pmb{a})\vee(s_2\wedge \pmb{c}),
  \]
  \[
    s'_1=(s_0\wedge \pmb{b})\vee(s_2\wedge \pmb{b}),
  \]
  \[
    s'_2=(s_1\wedge \pmb{c})\vee(s_2\wedge \pmb{a})\vee(s_2\wedge \pmb{c}).
  \]
  The initial configuration has $q_0=s_0=1$ and $s_1=s_2=0$, and the accept
  line is $s_2\wedge\text{\<EOS\>}$.

  \item \emph{DFA (connected portion only).}
  Starting from $\{q_0,s_0\}$, the reachable states are
  \[
    Q_0=\{q_0,s_0\},\qquad
    Q_1=\{s_0\},\qquad
    Q_2=\{s_1\},\qquad
    Q_3=\emptyset,\qquad
    Q_4=\{s_0,s_2\}.
  \]
  Their transitions are
  \[
  \begin{array}{c|ccc}
   & \pmb{a} & \pmb{b} & \pmb{c}\\\hline
  Q_0 & Q_1 & Q_2 & Q_1\\
  Q_1 & Q_1 & Q_2 & Q_1\\
  Q_2 & Q_3 & Q_3 & Q_4\\
  Q_3 & Q_3 & Q_3 & Q_3\\
  Q_4 & Q_4 & Q_2 & Q_4
  \end{array}
  \]
  The only final state is $Q_4$.
\end{enumerate}

%% 4.8
\item \textbf{Complementing the final states of an NDFA does not yield a machine
accepting the complement language.}

For a DFA, $L(\overline{A})=\Sigma^*\setminus L(A)$ because every word follows
exactly one path and is accepted iff it ends in a final state.  Flipping final
and nonfinal states reverses acceptance on that unique path.

For an NDFA, a word $w$ may follow \emph{multiple} paths, some ending in final
states and some not.  Under Definition~4.4, $w\in L(A)$ iff \emph{at least one}
path ends in a final state.  If we complement the final states, a word is
accepted by the modified machine iff at least one path ends in a \emph{formerly
nonfinal} state.  This is \emph{not} the same as $w\notin L(A)$, because a word
can simultaneously have a path to a (formerly) final state and a path to a
(formerly) nonfinal state; the complemented machine would accept it, but it was
already in $L(A)$.

\textbf{Example.}  Let $\Sigma=\{\pmb{a}\}$, $S=\{s_0,s_1\}$, $S_0=\{s_0\}$,
$F=\{s_1\}$, $\delta(s_0,\pmb{a})=\{s_0,s_1\}$, $\delta(s_1,\pmb{a})=\emptyset$.
Then $\pmb{a}\in L(A)$ because $s_1\in\delta(s_0,\pmb{a})$.  After complementing,
$F'=\{s_0\}$, and $\pmb{a}$ is still accepted (via the path $s_0\to s_0\in F'$),
yet $\pmb{a}\notin\Sigma^*\setminus L(A)$ since $\pmb{a}\in L(A)$.  Hence the
complement-of-final-states trick fails for NDFAs.

%% 4.9
\item \textbf{Given an NDFA $A$ without $\lambda$-transitions, construct a
$\lambda$-NDFA $A'$ with exactly one start state, exactly one final state, and
$L(A')=L(A)$.}

\textbf{Construction.}  Let $A=\langle\Sigma,S,S_0,\delta,F\rangle$.  Define
$A'=\langle\Sigma,S',\{q_0\},\delta',\{q_f\}\rangle$ where:
\begin{itemize}
  \item $S'=S\cup\{q_0,q_f\}$ (two new states).
  \item $q_0$ is the unique start state.
  \item $\delta'(q_0,\lambda)=S_0$ (a $\lambda$-transition from $q_0$ to every
        start state of $A$).
  \item For all $s\in S$ and $\pmb{a}\in\Sigma$: $\delta'(s,\pmb{a})=\delta(s,\pmb{a})$.
  \item $\delta'(f,\lambda)\ni q_f$ for every $f\in F$ ($\lambda$-transitions
        from each old final state to the new unique final state $q_f$).
  \item $\delta'(q_f,\pmb{a})=\emptyset$ for all $\pmb{a}$ (and no $\lambda$-transitions
        out of $q_f$).
\end{itemize}

\textbf{Correctness.}  A word $w\in\Sigma^+$ is accepted by $A$ iff there is a
path $q_0\xrightarrow{\lambda}s_i\xrightarrow{w}f$ for some $s_i\in S_0$ and
$f\in F$; by the $\lambda$-transition $f\xrightarrow{\lambda}q_f$, this same
path exists in $A'$ ending at $q_f$.  Conversely, any accepting path in $A'$
uses the initial $\lambda$-step to some $s_i\in S_0$, processes $w$ within $S$
(by identical transitions), then uses a final $\lambda$-step to $q_f$, which
corresponds to an accepting path in $A$.

If $\lambda\in L(A)$, add the extra $\lambda$-transition
$\delta'(q_0,\lambda)\ni q_f$.  Then $\lambda$ is accepted while $A'$ still has
exactly one start state ($q_0$) and exactly one final state ($q_f$).  Hence
$L(A')=L(A)$ in all cases.

%% 4.10
\item \textbf{Can part (ii) of Definition~4.8 be deduced from parts (i) and (iii)?}

Definition~4.8 defines $\DBAR$ for NDFAs:
\begin{enumerate}[(i)]
  \item $\DBAR(s,\lambda)=\{s\}$,
  \item $(\forall\pmb{a}\in\Sigma) \DBAR(s,\pmb{a})=\delta(s,\pmb{a})$,
  \item $(\forall x\in\Sigma^*)(\forall\pmb{a}\in\Sigma) \DBAR(s,x\pmb{a})
        =\bigcup_{t\in\DBAR(s,x)}\delta(t,\pmb{a})$.
\end{enumerate}
Part (ii) concerns single-letter strings.  Setting $x=\lambda$ in part (iii):
\[
  \DBAR(s,\lambda\pmb{a})
  =\bigcup_{t\in\DBAR(s,\lambda)}\delta(t,\pmb{a})
  =\bigcup_{t\in\{s\}}\delta(t,\pmb{a})
  =\delta(s,\pmb{a}),
\]
where we used part (i) to obtain $\DBAR(s,\lambda)=\{s\}$ and the fact that
$\lambda\pmb{a}=\pmb{a}$.  Since $\lambda\pmb{a}=\pmb{a}$, this gives
$\DBAR(s,\pmb{a})=\delta(s,\pmb{a})$, which is exactly part (ii).

\textbf{Conclusion.}  Yes, part (ii) can be derived from parts (i) and (iii).
It is included explicitly to make the base case of inductive proofs more
transparent, but it is logically redundant.

%% 4.11
\item \textbf{Alternative construction for removing $\lambda$-transitions.}

Define $S'=S$, $S'_0=\Lambda(S_0)$, $F'=F$, and
$\delta'(s,\pmb{a})=\DBAR_{\lambda}(s,\pmb{a})$ for all $s\in S$, $\pmb{a}\in\Sigma$.

\textbf{This construction does work.}  Let
\[
A'=\langle\Sigma,S,\Lambda(S_0),\delta',F\rangle.
\]
We show $L(A')=L(A_{\lambda})$.

First, for every $s\in S$ and every nonempty word $x\in\Sigma^+$,
\[
\overline{\delta'}(s,x)=\DBAR_{\lambda}(s,x).
\]
We prove this by induction on $|x|$.

\textbf{Base case.}  If $x=\pmb{a}\in\Sigma$, then by definition of $\delta'$,
\[
\overline{\delta'}(s,\pmb{a})=\delta'(s,\pmb{a})=\DBAR_{\lambda}(s,\pmb{a}).
\]

\textbf{Inductive step.}  Assume $\overline{\delta'}(s,x)=\DBAR_{\lambda}(s,x)$
for some $x\in\Sigma^+$.  Let $\pmb{a}\in\Sigma$.  Then
\begin{align*}
\overline{\delta'}(s,x\pmb{a})
&=\bigcup_{t\in\overline{\delta'}(s,x)}\delta'(t,\pmb{a}) \\
&=\bigcup_{t\in\DBAR_{\lambda}(s,x)}\DBAR_{\lambda}(t,\pmb{a})
\qquad\text{(induction hypothesis and definition of $\delta'$).}
\end{align*}
Now $\DBAR_{\lambda}(s,x)$ is $\lambda$-closed, because every value of
$\DBAR_{\lambda}$ is obtained by applying $\Lambda$.  Hence if
$t\in\DBAR_{\lambda}(s,x)$, then $\Lambda(t)\subseteq\DBAR_{\lambda}(s,x)$, and
so
\[
\DBAR_{\lambda}(t,\pmb{a})
=\Lambda\Bigl(\bigcup_{q\in\Lambda(t)}\delta_{\lambda}(q,\pmb{a})\Bigr)
\subseteq
\Lambda\Bigl(\bigcup_{q\in\DBAR_{\lambda}(s,x)}\delta_{\lambda}(q,\pmb{a})\Bigr)
=\DBAR_{\lambda}(s,x\pmb{a}).
\]
This shows
\[
\bigcup_{t\in\DBAR_{\lambda}(s,x)}\DBAR_{\lambda}(t,\pmb{a})
\subseteq \DBAR_{\lambda}(s,x\pmb{a}).
\]
For the reverse inclusion, if $r\in\DBAR_{\lambda}(s,x\pmb{a})$, then by
definition there is some $q\in\DBAR_{\lambda}(s,x)$ and some
$u\in\Lambda(q)\subseteq\DBAR_{\lambda}(s,x)$ such that $r$ is reached from $u$
by reading $\pmb{a}$ followed by $\lambda$-moves.  Therefore
$r\in\DBAR_{\lambda}(u,\pmb{a})$ for some $u\in\DBAR_{\lambda}(s,x)$, so
\[
\DBAR_{\lambda}(s,x\pmb{a})
\subseteq
\bigcup_{t\in\DBAR_{\lambda}(s,x)}\DBAR_{\lambda}(t,\pmb{a}).
\]
Thus
\[
\overline{\delta'}(s,x\pmb{a})=\DBAR_{\lambda}(s,x\pmb{a}),
\]
completing the induction.

\textbf{Case 1: $x=\lambda$.}  Then
\[
\lambda\in L(A')\iff \Lambda(S_0)\cap F\neq\emptyset.
\]
But this says exactly that some start state of $A_{\lambda}$ can reach a final
state by $\lambda$-moves alone, which is the condition $\lambda\in L(A_{\lambda})$.

\textbf{Case 2: $x\in\Sigma^+$.}  Using the identity above,
\[
x\in L(A')\iff
\Bigl(\bigcup_{t\in\Lambda(S_0)}\overline{\delta'}(t,x)\Bigr)\cap F\neq\emptyset
\iff
\Bigl(\bigcup_{t\in\Lambda(S_0)}\DBAR_{\lambda}(t,x)\Bigr)\cap F\neq\emptyset.
\]
If $t\in\Lambda(S_0)$, then $t\in\Lambda(s_0)$ for some $s_0\in S_0$, so the
initial $\lambda$-moves taking $s_0$ to $t$ may be prefixed to any path counted
in $\DBAR_{\lambda}(t,x)$. Hence
\[
\bigcup_{t\in\Lambda(S_0)}\DBAR_{\lambda}(t,x)
=
\bigcup_{s_0\in S_0}\DBAR_{\lambda}(s_0,x),
\]
and therefore
\[
x\in L(A')\iff
\Bigl(\bigcup_{s_0\in S_0}\DBAR_{\lambda}(s_0,x)\Bigr)\cap F\neq\emptyset
\iff x\in L(A_{\lambda}).
\]
So this is a valid alternative to Definition~4.9: instead of adding $q_0$, it
enlarges the start-state set to $\Lambda(S_0)$.

%% 4.12
\item \textbf{Algorithm for union of two FAD languages using $\lambda$-NDFAs.}

Let $L_1=L(A_1)$ and $L_2=L(A_2)$ where $A_1$ and $A_2$ are $\lambda$-NDFAs.
By Exercise~4.9 we may assume each $A_i$ has exactly one start state $q^i_0$ and
exactly one final state $q^i_f$.  Construct $A'$ as follows:
\begin{enumerate}[(1)]
  \item Introduce a new start state $q_0$ and a new final state $q_f$.
  \item Add $\lambda$-transitions: $q_0\xrightarrow{\lambda}q^1_0$ and
        $q_0\xrightarrow{\lambda}q^2_0$.
  \item Add $\lambda$-transitions: $q^1_f\xrightarrow{\lambda}q_f$ and
        $q^2_f\xrightarrow{\lambda}q_f$.
  \item Keep all existing transitions of $A_1$ and $A_2$.
\end{enumerate}
Then $L(A')=L(A_1)\cup L(A_2)$: a word is accepted iff it is accepted by $A_1$
or by $A_2$, because the $\lambda$-transitions route the computation into one
of the two machines and collect the result at $q_f$.

%% 4.13
\item \textbf{Examples of three NDFAs such that the corresponding $A^d$ is
respectively not connected, not reduced, and minimal.}

\textbf{(a) $A^d$ is not connected.}  Any NDFA in which some subsets of states
are never reachable from the start set under the transition function will
produce an $A^d$ with unreachable states.  For example, let
$\Sigma=\{\pmb{a}\}$, $S=\{s_0,s_1\}$, $S_0=\{s_0\}$, $F=\{s_1\}$,
$\delta(s_0,\pmb{a})=\{s_1\}$, $\delta(s_1,\pmb{a})=\emptyset$.
Then $A^d$ has states $\{\emptyset,\{s_0\},\{s_1\},\{s_0,s_1\}\}$ but
$\{s_0,s_1\}$ and $\{s_1\}$ with $\delta^d(\{s_1\},\pmb{a})=\emptyset$ are
reachable, while $\{s_0,s_1\}$ is not reached from $\{s_0\}$.

\textbf{(b) $A^d$ is not reduced.}  Use a different example:
\[
A=\langle\{\pmb{a},\pmb{b}\},\{p,u,v\},\{p\},\delta,\{u,v\}\rangle
\]
where
\[
\delta(p,\pmb{a})=\{u\},\qquad \delta(p,\pmb{b})=\{v\},
\]
and all other transitions are empty.  Then the reachable subset states of $A^d$
are
\[
\{p\},\qquad \{u\},\qquad \{v\},\qquad \emptyset.
\]
Both $\{u\}$ and $\{v\}$ are final, and from either one every input leads to
$\emptyset$. Therefore
\[
\{u\}\ E_{A^d}\ \{v\},
\]
so $A^d$ is not reduced.

\textbf{(c) $A^d$ is minimal.}  Here we use a different example.  Let
$A=\langle\{\pmb{a},\pmb{b}\},\{s\},\{s\},\delta,\{s\}\rangle$ where
\[
\delta(s,\pmb{a})=\{s\},\qquad \delta(s,\pmb{b})=\emptyset.
\]
Then $L(A)=\pmb{a}^*$.  The deterministic machine $A^d$ has exactly two states,
$\{s\}$ and $\emptyset$, with
\[
\delta^d(\{s\},\pmb{a})=\{s\},\quad \delta^d(\{s\},\pmb{b})=\emptyset,\quad
\delta^d(\emptyset,\pmb{a})=\delta^d(\emptyset,\pmb{b})=\emptyset.
\]
This is the standard 2-state DFA for $\pmb{a}^*$ over the alphabet
$\{\pmb{a},\pmb{b}\}$, so $A^d$ is minimal.

%% 4.14
\item \textbf{Why was the extra state $q_0$ included in $A^{\sigma}_\lambda$?}

The role of $q_0$ is to preserve acceptance of $\lambda$ while leaving the
original start states and all ordinary transitions unchanged.  In
Definition~4.9 the extra state $q_0$ is disconnected from the rest of the
machine, and
\[
q_0\in F^{\sigma}\iff \lambda\in L(A_{\lambda}).
\]
Thus $q_0$ serves as a separate witness for the empty word.

\textbf{Example.}  Let
\[
A_{\lambda}=\langle\{\pmb{a}\},\{s_0,f\},\{s_0\},\delta_{\lambda},\{f\}\rangle
\]
where $\delta_{\lambda}(s_0,\lambda)=\{f\}$ and all ordinary transitions are
empty.  Then $\lambda\in L(A_{\lambda})$ because $f\in\Lambda(s_0)$.

If we try to remove $\lambda$-moves \emph{without} adding $q_0$ and keep the
same start state $s_0$, then the new machine has no way to accept $\lambda$:
no input is read, and $s_0$ is not final.  The disconnected state $q_0$ fixes
exactly this problem.  We keep $s_0$ for all nonempty words, and we make $q_0$
final so that the empty word is still accepted.

%% 4.15
\item \textbf{Union construction using NDFAs \emph{without} $\lambda$-transitions.}

\begin{enumerate}[(a)]
  \item \textbf{Algorithm.}  Let $A_1=\langle\Sigma,S_1,S^1_0,\delta_1,F_1\rangle$
  and $A_2=\langle\Sigma,S_2,S^2_0,\delta_2,F_2\rangle$ (assume $S_1\cap S_2=\emptyset$).
  Construct $A'=\langle\Sigma,S_1\cup S_2,S^1_0\cup S^2_0,\delta',F_1\cup F_2\rangle$
  where $\delta'(s,\pmb{a})=\delta_i(s,\pmb{a})$ for $s\in S_i$.
  The machine $A'$ starts in the combined set of initial states and accepts
  any word accepted by either $A_1$ or $A_2$.  Thus $L(A')=L(A_1)\cup L(A_2)$.

  \item \textbf{Comparison.}  This is actually \emph{easier} than the
  $\lambda$-transition approach (Exercise~4.12): no new states are needed at
  all, since we simply merge the state sets and initial-state sets.

  \item \textbf{Can we ensure exactly one start state and one final state?}
  Not in full generality.  The disjoint-union construction in part~(a) gives a
  correct union machine, but it usually has several start states and several
  final states.  If the resulting language does \emph{not} contain $\lambda$,
  then Exercise~4.22 shows how to convert a no-$\lambda$ NDFA into another
  no-$\lambda$ NDFA with exactly one start state and exactly one final state.
  But Exercise~4.23 shows that this can fail when $\lambda$ is in the language.
  So the correct conclusion is: one start state and one final state can be
  forced under the hypothesis $\lambda\notin L$, but not uniformly in general.
\end{enumerate}

%% 4.16
\item Consider the $\lambda$-NDFA $A^{\sigma}_\lambda$ from Example~4.15 (the
machine without $\lambda$-transitions corresponding to the automaton of
Example~4.14).

\begin{enumerate}[(a)]
  \item \emph{Circuit diagram (no SOS or EOS).}
  The equivalent $\lambda$-free machine has the five states
  $q_0,s_0,s_1,s_2,s_3$, so we use one activity bit for each of them.  The
  next-state equations are
  \[
    q'_0=0,\qquad s'_0=0,
  \]
  \[
    s'_1=(s_0\wedge \pmb{a})\vee(s_3\wedge \pmb{b}),
  \]
  \[
    s'_2=(s_0\wedge \pmb{a})\vee(s_0\wedge \pmb{c})\vee(s_0\wedge \pmb{d})
       \vee(s_1\wedge \pmb{c})\vee(s_1\wedge \pmb{d})
       \vee(s_3\wedge \pmb{b})\vee(s_2\wedge \pmb{d}),
  \]
  \[
    s'_3=(s_0\wedge \pmb{b})\vee(s_1\wedge \pmb{b}).
  \]
  The initial configuration has $q_0=s_0=1$ and $s_1=s_2=s_3=0$.  Since
  $q_0$ and $s_2$ are final states in Example~4.15, the accept output is
  $q_0\vee s_2$.

  \item \emph{Convert to $A^{\sigma d}_\lambda$ (connected portion only).}
  In Example~4.15, the start states of $A^{\sigma}_\lambda$ are
  \[
  \{q_0,s_0\},
  \]
  and the final states are $q_0$ and $s_2$.  The connected part of the
  deterministic equivalent has the reachable states
  \[
  R_0=\{q_0,s_0\},\quad
  R_1=\{s_1,s_2\},\quad
  R_2=\{s_3\},\quad
  R_3=\{s_2\},\quad
  R_4=\emptyset.
  \]
  Their transitions are
  \[
  \begin{array}{c|cccc}
   & \pmb{a} & \pmb{b} & \pmb{c} & \pmb{d}\\\hline
  R_0 & R_1 & R_2 & R_3 & R_3\\
  R_1 & R_4 & R_2 & R_3 & R_3\\
  R_2 & R_4 & R_1 & R_4 & R_4\\
  R_3 & R_4 & R_4 & R_4 & R_3\\
  R_4 & R_4 & R_4 & R_4 & R_4
  \end{array}
  \]
  The final states are
  \[
  R_0,\ R_1,\ R_3,
  \]
  because these are exactly the reachable subsets containing $q_0$ or $s_2$.
\end{enumerate}

%% 4.17
\item \begin{enumerate}[(a)]
  \item \textbf{Proof that $\DBAR(s,\pmb{a})=\delta(s,\pmb{a})$ for any NDFA
  without $\lambda$-transitions.}

  By Definition~4.8(ii), $\DBAR(s,\pmb{a})=\delta(s,\pmb{a})$ for all
  $s\in S$ and $\pmb{a}\in\Sigma$.  This is immediate from the definition;
  alternatively, apply part (iii) with $x=\lambda$:
  \[
    \DBAR(s,\lambda\pmb{a})
    =\bigcup_{t\in\DBAR(s,\lambda)}\delta(t,\pmb{a})
    =\bigcup_{t\in\{s\}}\delta(t,\pmb{a})
    =\delta(s,\pmb{a}),
  \]
  using part (i): $\DBAR(s,\lambda)=\{s\}$.  Since $\lambda\pmb{a}=\pmb{a}$,
  we conclude $\DBAR(s,\pmb{a})=\delta(s,\pmb{a})$. \qed

  \item \textbf{Counterexample for $\lambda$-NDFAs.}

  Let $\Sigma=\{\pmb{a}\}$, $S=\{s_0,s_1\}$, $S_0=\{s_0\}$, $F=\{s_1\}$,
  $\delta_\lambda(s_0,\lambda)=\{s_1\}$, $\delta_\lambda(s_0,\pmb{a})=\emptyset$,
  $\delta_\lambda(s_1,\pmb{a})=\emptyset$.

  Then $\Lambda(s_0)=\{s_0,s_1\}$.  So
  \[
    \DBAR_\lambda(s_0,\pmb{a})
    =\Lambda \Bigl(\bigcup_{t\in\Lambda(s_0)}\delta_\lambda(t,\pmb{a})\Bigr)
    =\Lambda(\emptyset\cup\emptyset)=\emptyset,
  \]
  but $\delta_\lambda(s_0,\pmb{a})=\emptyset$ also.  Let us try a different
  example: $\delta_\lambda(s_0,\lambda)=\{s_1\}$, $\delta_\lambda(s_1,\pmb{a})=\{s_1\}$.
  Then $\delta_\lambda(s_0,\pmb{a})=\emptyset$, yet
  \[
    \DBAR_\lambda(s_0,\pmb{a})
    =\Lambda \Bigl(\bigcup_{t\in\{s_0,s_1\}}\delta_\lambda(t,\pmb{a})\Bigr)
    =\Lambda(\{s_1\})=\{s_1\}\neq\emptyset=\delta_\lambda(s_0,\pmb{a}).
  \]
  This demonstrates that $\DBAR_\lambda(s,\pmb{a})\neq\delta_\lambda(s,\pmb{a})$
  is possible for $\lambda$-NDFAs.
\end{enumerate}

%% 4.18
\item Consider the NDFA accepting the original language $L$ from Example~4.10:
$L=\{x\in\{\pmb{a},\pmb{b}\}^*\mid x\text{ begins with }\pmb{a}\vee x\text{ contains }\pmb{ba}\}$.

\begin{enumerate}[(a)]
  \item \emph{Circuit diagram (no SOS/EOS).}  The NDFA has a start state
  $q_0$, two final states $q_1,q_3$, and the transitions of Figure~4.10.
  Using one activity bit per state, the next-state equations are
  \[
    q'_0=0,\qquad
    q'_1=(q_0\wedge \pmb{a})\vee(q_1\wedge \pmb{a})\vee(q_1\wedge \pmb{b}),
  \]
  \[
    q'_2=(q_0\wedge \pmb{b})\vee(q_2\wedge \pmb{b}),\qquad
    q'_3=(q_2\wedge \pmb{a})\vee(q_3\wedge \pmb{a})\vee(q_3\wedge \pmb{b}).
  \]
  Initially only $q_0$ is active, and the accept line is $q_1\vee q_3$.

  \item \emph{DFA (connected portion only).}  Apply the subset construction
  to Figure~4.10.  Starting from $\{q_0\}$, the reachable states are
  \[
    Q_0=\{q_0\},\qquad Q_1=\{q_1\},\qquad Q_2=\{q_2\},\qquad Q_3=\{q_3\}.
  \]
  Their transitions are
  \[
  \begin{array}{c|cc}
   & \pmb{a} & \pmb{b}\\\hline
  Q_0 & Q_1 & Q_2\\
  Q_1 & Q_1 & Q_1\\
  Q_2 & Q_3 & Q_2\\
  Q_3 & Q_3 & Q_3
  \end{array}
  \]
  The final states are $Q_1$ and $Q_3$.
\end{enumerate}

%% 4.19
\item Consider the NDFA accepting $L^r=\{x\in\{\pmb{a},\pmb{b}\}^*\mid x\text{ ends with }\pmb{a}\vee x\text{ contains }\pmb{ab}\}$ from Example~4.10.

\begin{enumerate}[(a)]
  \item \emph{Circuit diagram (no SOS/EOS).}  This NDFA is obtained by
  reversing the arrows in Figure~4.10 and swapping the start and final
  states.  Using activity bits $q_0,q_1,q_2,q_3$, the next-state equations are
  \[
    q'_0=(q_1\wedge \pmb{a})\vee(q_2\wedge \pmb{b}),
  \]
  \[
    q'_1=(q_1\wedge \pmb{a})\vee(q_1\wedge \pmb{b}),\qquad
    q'_2=(q_3\wedge \pmb{a})\vee(q_2\wedge \pmb{b}),
  \]
  \[
    q'_3=(q_3\wedge \pmb{a})\vee(q_3\wedge \pmb{b}).
  \]
  The initial configuration has $q_1=q_3=1$ and $q_0=q_2=0$, and the accept
  line is $q_0$.

  \item \emph{DFA (connected portion only).}  From Figure~4.11, the NDFA has
  start states $\{q_1,q_3\}$ and only one final state, $q_0$.  Its nonempty
  transitions are
  \[
  \delta(q_1,\pmb{a})=\{q_1,q_0\},\qquad \delta(q_1,\pmb{b})=\{q_1\},
  \]
  \[
  \delta(q_2,\pmb{b})=\{q_2,q_0\},\qquad
  \delta(q_3,\pmb{a})=\{q_3,q_2\},\qquad
  \delta(q_3,\pmb{b})=\{q_3\}.
  \]
  Starting from $Q_0=\{q_1,q_3\}$, Definition~4.5 gives
  \[
  \delta^d(Q_0,\pmb{a})=\{q_0,q_1,q_2,q_3\}=Q_1,\qquad
  \delta^d(Q_0,\pmb{b})=\{q_1,q_3\}=Q_0.
  \]
  From $Q_1$ we get
  \[
  \delta^d(Q_1,\pmb{a})=Q_1,\qquad \delta^d(Q_1,\pmb{b})=Q_1.
  \]
  Therefore the connected DFA has just two states:
  \[
  Q_0=\{q_1,q_3\}\quad\text{(nonfinal)},\qquad
  Q_1=\{q_0,q_1,q_2,q_3\}\quad\text{(final)},
  \]
  with transition table
  \[
  \begin{array}{c|cc}
   & \pmb{a} & \pmb{b}\\\hline
  Q_0 & Q_1 & Q_0\\
  Q_1 & Q_1 & Q_1
  \end{array}
  \]
\end{enumerate}

%% 4.20
\item \textbf{Are the pumping lemma arguments still valid for NDFAs?}

Not as \emph{machine} arguments about an NDFA.  In Chapter~2, the pumping lemma
was proved for DFAs by following the \emph{single} path determined by a long
input and finding a repeated state on that path.  An NDFA does not behave that
way, because one word may have many possible paths.

However, the pumping lemma conclusion \emph{does} still hold for every language
accepted by an NDFA.  If $A$ is an NDFA, then by Theorem~4.1 there is an
equivalent DFA $A^d$ with $L(A)=L(A^d)$.  Since $L(A^d)$ is FAD, the Chapter~2
pumping lemma applies to that language.  Therefore any language accepted by an
NDFA satisfies the pumping lemma, but the proof is obtained by passing to the
equivalent DFA, not by reusing the Chapter~2 path argument directly on the NDFA.

%% 4.21
\item \textbf{Does the conclusion of Theorem~2.7 still hold for NDFAs?}

Yes, but again this is a statement about the \emph{language}, not about the
internal states of an arbitrary NDFA.  Theorem~2.7 says that a language $L$ is
FAD iff $\equiv_L$ has finitely many equivalence classes.  By Theorem~4.1, a
language is accepted by some NDFA iff it is accepted by some DFA.  Therefore
the conclusion of Theorem~2.7 still holds for languages accepted by NDFAs:
\[
L \text{ is accepted by an NDFA }
\Leftrightarrow
L \text{ is FAD }
\Leftrightarrow
\equiv_L \text{ has finitely many classes.}
\]
What does \emph{not} follow is that the states of a given NDFA correspond to the
equivalence classes of $\equiv_L$; that correspondence belongs to the minimal
DFA, not to an arbitrary nondeterministic machine.

%% 4.22
\item \textbf{Later no-$\lambda$ exercise: given an NDFA $A$ without $\lambda$-transitions
with $\lambda\notin L(A)$, construct $A''$ with exactly one start state, one
final state, and $L(A'')=L(A)$.}

\textbf{Construction.}  Let $A=\langle\Sigma,S,S_0,\delta,F\rangle$ with
$\lambda\notin L(A)$, so $S_0\cap F=\emptyset$.  Define
\[
A''=\langle \Sigma,  S\cup\{q_0,q_f\},  \{q_0\},  \delta'',  \{q_f\}\rangle
\]
where $q_0$ and $q_f$ are new states and:
\begin{enumerate}[1.]
  \item For each $\pmb{a}\in\Sigma$, let
  \[
  \delta''(q_0,\pmb{a})=
  \Bigl(\bigcup_{s\in S_0}\delta(s,\pmb{a})\Bigr)\cup
  \left\{\begin{array}{ll}
  \{q_f\} & \text{if }(\bigcup_{s\in S_0}\delta(s,\pmb{a}))\cap F\neq\emptyset,\\
  \emptyset & \text{otherwise.}
  \end{array}\right.
  \]
  \item For each old state $s\in S$ and each $\pmb{a}\in\Sigma$, let
  \[
  \delta''(s,\pmb{a})=
  \delta(s,\pmb{a})\cup
  \left\{\begin{array}{ll}
  \{q_f\} & \text{if }\delta(s,\pmb{a})\cap F\neq\emptyset,\\
  \emptyset & \text{otherwise.}
  \end{array}\right.
  \]
  \item $\delta''(q_f,\pmb{a})=\emptyset$ for every $\pmb{a}\in\Sigma$.
\end{enumerate}

This machine has exactly one start state ($q_0$) and exactly one final state
($q_f$).

\textbf{Why this preserves the language.}  Since $\lambda\notin L(A)$, every
accepted word has positive length.  Suppose
$w=\pmb{a}_1\cdots \pmb{a}_n\in L(A)$.  Then there is a path in $A$ from some
start state to some final state that reads all of $w$.  In $A''$, use the same
path for the first $n-1$ letters.  On the last letter, when the old path would
move into a final state, choose the extra target $q_f$ instead.  Thus
$w\in L(A'')$.

Conversely, any accepting path in $A''$ must end in $q_f$.  The only way to
enter $q_f$ is on some actual input symbol $\pmb{a}$ from $q_0$ or from an old
state $s$, and this happens only when the corresponding old transition set
contains at least one state of $F$.  Replacing that last step by a move into
one of those original final states yields an accepting path in $A$.  Hence
$L(A'')=L(A)$.

%% 4.23
\item \textbf{If $\lambda\in L(A)$, the three conditions of Exercise~4.22
may be incompatible (without $\lambda$-transitions).}

\textbf{Example.}  Let
\[
A=\langle\{\pmb{a}\},\{s_0,s_1\},\{s_0\},\delta,\{s_0,s_1\}\rangle
\]
where $\delta(s_0,\pmb{a})=\{s_1\}$ and $\delta(s_1,\pmb{a})=\emptyset$.  Then
\[
L(A)=\{\lambda,\pmb{a}\}.
\]

We claim there is no NDFA without $\lambda$-transitions that has exactly one
start state, exactly one final state, and also accepts exactly
$\{\lambda,\pmb{a}\}$.

Assume to the contrary that
\[
B=\langle\{\pmb{a}\},S,\{q\},\delta_B,\{f\}\rangle
\]
has those three properties and $L(B)=\{\lambda,\pmb{a}\}$.  Since $\lambda\in
L(B)$, the unique start state must also be the unique final state:
\[
q=f.
\]
Since $\pmb{a}\in L(B)$ and there are no $\lambda$-transitions, there must be a
transition on $\pmb{a}$ from $q$ to the only final state, namely to $q$
itself.  Thus $q\in\delta_B(q,\pmb{a})$.

But then the same move can be used repeatedly, so every positive power of
$\pmb{a}$ is accepted:
\[
\pmb{a}^n\in L(B)\qquad\text{for all }n\geq 1.
\]
In particular, $\pmb{aa}\in L(B)$, contradicting $L(B)=\{\lambda,\pmb{a}\}$.
Therefore the three conditions may indeed be incompatible.

%% 4.24
\item \textbf{Prove Lemma~4.2.}

\textbf{Lemma~4.2:} For any $s\in S$ and any $x\in\Sigma^+$,
$\overline{\delta^{\sigma}_\lambda}(s,x)=\DBAR_\lambda(s,x)$.

\begin{proof}
We proceed by induction on $|x|=k\ge 1$.

\textbf{Base case ($k=1$).}  Let $x=\pmb{a}\in\Sigma$.  By definition of
$\delta^{\sigma}_\lambda$:
\[
  \delta^{\sigma}_\lambda(s,\pmb{a})
  =\Lambda \Bigl(\bigcup_{t\in\Lambda(s)}\delta_\lambda(t,\pmb{a})\Bigr)
  =\DBAR_\lambda(s,\pmb{a}).
\]
So $\overline{\delta^{\sigma}_\lambda}(s,\pmb{a})=\delta^{\sigma}_\lambda(s,\pmb{a})=\DBAR_\lambda(s,\pmb{a})$.

\textbf{Inductive step.}  Assume the result holds for all words of length $k$.
Let $y=x\pmb{a}$ with $|x|=k$ and $\pmb{a}\in\Sigma$.  Then:
\begin{align*}
  \overline{\delta^{\sigma}_\lambda}(s,x\pmb{a})
  &= \bigcup_{t\in\overline{\delta^{\sigma}_\lambda}(s,x)}
      \delta^{\sigma}_\lambda(t,\pmb{a}) \\
  &= \bigcup_{t\in\DBAR_\lambda(s,x)}
      \delta^{\sigma}_\lambda(t,\pmb{a}) \quad\text{(induction hypothesis)} \\
  &= \bigcup_{t\in\DBAR_\lambda(s,x)}
      \DBAR_\lambda(t,\pmb{a}) \quad\text{(base case)}.
\end{align*}
Let $D=\DBAR_\lambda(s,x)$.  Since every value of $\DBAR_\lambda$ is obtained
by applying $\Lambda$, the set $D$ is $\lambda$-closed.  Therefore
\begin{align*}
\bigcup_{t\in D}\DBAR_\lambda(t,\pmb{a})
&=
\bigcup_{t\in D}
\Lambda\Bigl(\bigcup_{q\in\Lambda(t)}\delta_\lambda(q,\pmb{a})\Bigr) \\
&=
\Lambda\Bigl(\bigcup_{t\in D}\bigcup_{q\in\Lambda(t)}\delta_\lambda(q,\pmb{a})\Bigr) \\
&=
\Lambda\Bigl(\bigcup_{q\in D}\delta_\lambda(q,\pmb{a})\Bigr)
\qquad\text{(because $\Lambda(t)\subseteq D$ for every $t\in D$)} \\
&= \DBAR_\lambda(s,x\pmb{a}).
\end{align*}
By induction, the result holds for all $x\in\Sigma^+$. \qed
\end{proof}

%% 4.25
\item \textbf{A DFA can be viewed as an NDFA, and $L(A^n)=L(A)$.}

\textbf{Construction.}  Given a DFA $A=\langle\Sigma,S,s_0,\delta,F\rangle$
(with $\delta:S\times\Sigma\to S$ a total function), define the NDFA
$A^n=\langle\Sigma,S,\{s_0\},\delta^n,F\rangle$ where
$\delta^n:S\times\Sigma\to\wp(S)$ is given by $\delta^n(s,\pmb{a})=\{\delta(s,\pmb{a})\}$.

$A^n$ is a valid NDFA: $S$ is finite, $S_0^n=\{s_0\}$ is a nonempty set of
start states, and $\delta^n$ maps to subsets of $S$.

\textbf{Claim: $L(A^n)=L(A)$.}

\begin{proof}
By induction on $|x|$, we show $\DBAR^n(s_0,x)=\{\deltabar{}(s_0,x)\}$ (i.e., the
nondeterministic extended transition function applied to the singleton start
always yields the singleton $\{\deltabar{}(s_0,x)\}$).

\emph{Base:} $\DBAR^n(s_0,\lambda)=\{s_0\}=\{\deltabar{}(s_0,\lambda)\}$.

\emph{Step:} $\DBAR^n(s_0,x\pmb{a})
=\bigcup_{t\in\DBAR^n(s_0,x)}\delta^n(t,\pmb{a})
=\bigcup_{t\in\{\deltabar{}(s_0,x)\}}\{\delta(t,\pmb{a})\}
=\{\delta(\deltabar{}(s_0,x),\pmb{a})\}
=\{\deltabar{}(s_0,x\pmb{a})\}$.

Hence $x\in L(A^n)$ iff $\DBAR^n(s_0,x)\cap F\neq\emptyset$ iff
$\deltabar{}(s_0,x)\in F$ iff $x\in L(A)$. \qed
\end{proof}

%% 4.26
\item \textbf{Show that $\overline{\delta^{\sigma}_\lambda}(s,\lambda)\neq\DBAR_\lambda(s,\lambda)$ in general.}

By definition:
\begin{itemize}
  \item $\overline{\delta^{\sigma}_\lambda}(s,\lambda)=\{s\}$ (definition of
    the extended transition function for zero-length input).
  \item $\DBAR_\lambda(s,\lambda)=\Lambda(s)$ (definition of $\DBAR_\lambda$ for $\lambda$).
\end{itemize}

These differ whenever $\Lambda(s)\neq\{s\}$, i.e., whenever there exists at
least one $\lambda$-transition out of $s$.

\textbf{Example.}  In Example~4.14, $\Lambda(s_0)=\{s_0,s_1,s_2\}$ but
$\overline{\delta^{\sigma}_\lambda}(s_0,\lambda)=\{s_0\}$.  Hence
$\overline{\delta^{\sigma}_\lambda}(s_0,\lambda)\neq\DBAR_\lambda(s_0,\lambda)$.

This is why Lemma~4.2 is stated only for $x\in\Sigma^+$ (words of length
$\ge 1$) and cannot be extended to $\lambda$.

%% 4.27
\item \textbf{Prove: $(\forall A,B\in\wp(S))(\forall\pmb{a}\in\Sigma) 
\delta^d(A\cup B,\pmb{a})=\delta^d(A,\pmb{a})\cup\delta^d(B,\pmb{a})$.}

\begin{proof}
By Definition~4.5, $\delta^d(Q,\pmb{a})=\bigcup_{q\in Q}\delta(q,\pmb{a})$
for any $Q\in\wp(S)$.  Therefore:
\begin{align*}
  \delta^d(A\cup B,\pmb{a})
  &= \bigcup_{q\in A\cup B}\delta(q,\pmb{a}) \\
  &= \Bigl(\bigcup_{q\in A}\delta(q,\pmb{a})\Bigr)\cup\Bigl(\bigcup_{q\in B}\delta(q,\pmb{a})\Bigr) \\
  &= \delta^d(A,\pmb{a})\cup\delta^d(B,\pmb{a}). \qed
\end{align*}
\end{proof}

%% 4.28
\item \textbf{Acceptance requiring \emph{every} reachable state to be final.}

\begin{enumerate}[(a)]
  \item \textbf{New acceptance definition.}  Define
  \[
  L_{\forall}(A)=\Bigl\{w\in\Sigma^* \,\Bigm|\,
  \bigcup_{q\in S_0}\DBAR(q,w)\neq\emptyset
  \ \wedge\
  \bigcup_{q\in S_0}\DBAR(q,w)\subseteq F\Bigr\}.
  \]
  That is, $w$ is accepted iff (i) at least one state is reachable and
  (ii) \emph{every} reachable state is final.

  \item \textbf{Modified $A^d$.}  Define
  \[
    A^{\forall}=\langle \Sigma,\wp(S),S_0,\delta^d,F^{\forall}\rangle,
  \]
  where
  \[
    \delta^d(Q,\pmb{a})=\bigcup_{q\in Q}\delta(q,\pmb{a})
  \]
  for each $Q\subseteq S$ and $\pmb{a}\in\Sigma$, and where the final-state set is
  \[
    F^{\forall}=\{Q\in\wp(S)\mid Q\neq\emptyset\wedge Q\subseteq F\}.
  \]
  That is, $Q$ is a final state of $A^{\forall}$ iff $Q$ is nonempty and
  all states in $Q$ are final states of $A$.

  For any word $w$, the subset construction yields
  \[
    \DBAR^d(S_0,w)=\bigcup_{q\in S_0}\DBAR(q,w).
  \]
  Therefore
  \[
    w\in L(A^{\forall})
    \Leftrightarrow
    \DBAR^d(S_0,w)\in F^{\forall}
  \]
  \[
    \Leftrightarrow
    \bigcup_{q\in S_0}\DBAR(q,w)\neq\emptyset
    \ \wedge\
    \bigcup_{q\in S_0}\DBAR(q,w)\subseteq F
    \Leftrightarrow
    w\in L_{\forall}(A).
  \]
  Hence $L_{\forall}(A)=L(A^{\forall})$.

  Note: this ``universal'' acceptance is still FAD (the modified $A^{\forall}$
  is a DFA), so $L_{\forall}(A)$ is always a FAD language.
\end{enumerate}

%% 4.29
\item \textbf{Draw the connected part of $T^d$, the deterministic equivalent of
the NDFA $T$ in Example~4.7.}

Let
\[
Q_0=\{s_0\},\quad
Q_1=\{s_0,s_1\},\quad
Q_2=\{s_0,s_2\},\quad
Q_3=\{s_0,s_1,s_3\},
\]
\[
Q_4=\{s_0,s_1,s_4\},\quad
Q_5=\{s_0,s_2,s_5\},\quad
Q_6=\{s_0,s_1,s_3,s_6\},
\]
\[
Q_7=\{s_0,s_1,s_4,s_6\},\quad
Q_8=\{s_0,s_2,s_6\},\quad
Q_9=\{s_0,s_1,s_6\},
\]
\[
Q_{10}=\{s_0,s_2,s_5,s_6\},\qquad
Q_{11}=\{s_0,s_6\}.
\]
These are exactly the reachable subsets from $\{s_0\}$ under Definition~4.5.
The accepting states are those containing $s_6$, namely
\[
Q_6,Q_7,Q_8,Q_9,Q_{10},Q_{11}.
\]
Their transitions are:
\[
\begin{array}{c|cc}
 & \pmb{0} & \pmb{1}\\\hline
Q_0 & Q_1 & Q_0\\
Q_1 & Q_1 & Q_2\\
Q_2 & Q_3 & Q_0\\
Q_3 & Q_4 & Q_2\\
Q_4 & Q_1 & Q_5\\
Q_5 & Q_6 & Q_0\\
Q_6 & Q_7 & Q_8\\
Q_7 & Q_9 & Q_{10}\\
Q_8 & Q_6 & Q_{11}\\
Q_9 & Q_9 & Q_8\\
Q_{10} & Q_6 & Q_{11}\\
Q_{11} & Q_9 & Q_{11}
\end{array}
\]
This is the connected part of $T^d$.

%% 4.30
\item \textbf{Modify $T$ to revert to a nonfinal state when $\pmb{000111}$ (EOT) is detected.}

Keep the original start-of-transmission arm for $\pmb{010010}$, but replace the
single final sink $s_6$ by six final ``recording'' states $r_0,\dots,r_5$ that
track the longest suffix of the data seen so far that is also a prefix of the
EOT string $\pmb{000111}$.  The only nonfinal states are the original
$s_0,\dots,s_5$.

\[
\begin{array}{c|cc}
 & \pmb{0} & \pmb{1}\\\hline
s_0 & \{s_0,s_1\} & \{s_0\}\\
s_1 & \emptyset & \{s_2\}\\
s_2 & \{s_3\} & \emptyset\\
s_3 & \{s_4\} & \emptyset\\
s_4 & \emptyset & \{s_5\}\\
s_5 & \{r_0\} & \emptyset\\
r_0 & \{r_1\} & \{r_0\}\\
r_1 & \{r_2\} & \{r_0\}\\
r_2 & \{r_3\} & \{r_0\}\\
r_3 & \{r_3\} & \{r_4\}\\
r_4 & \{r_1\} & \{r_5\}\\
r_5 & \{r_1\} & \{s_0\}
\end{array}
\]

Here $r_i$ means that recording is on and the last $i$ scanned bits match the
first $i$ bits of $\pmb{000111}$.  Thus $r_5\xrightarrow{\pmb{1}}s_0$ turns the
recorder off exactly when the full EOT signal $\pmb{000111}$ has just been
seen, while every other input keeps the machine in one of the final recording
states.

%% 4.31
\item Consider the NDFA $A$ from Example~4.14.

\begin{enumerate}[(a)]
  \item \textbf{Diagram of $A^{\sigma}_\lambda$.}
  For Example~4.14, $\lambda\in L(A_\lambda)$, so Definition~4.9 yields
  \[
    A^{\sigma}_\lambda=
    \langle\{\pmb{a},\pmb{b},\pmb{c},\pmb{d}\},
    \{q_0,s_0,s_1,s_2,s_3\},\{q_0,s_0\},\delta^{\sigma}_\lambda,\{q_0,s_2\}\rangle
  \]
  with transition table
  \[
  \begin{array}{c|cccc}
   & \pmb{a} & \pmb{b} & \pmb{c} & \pmb{d}\\\hline
  q_0 & \emptyset & \emptyset & \emptyset & \emptyset\\
  s_0 & \{s_1,s_2\} & \{s_3\} & \{s_2\} & \{s_2\}\\
  s_1 & \emptyset & \{s_3\} & \{s_2\} & \{s_2\}\\
  s_2 & \emptyset & \emptyset & \emptyset & \{s_2\}\\
  s_3 & \emptyset & \{s_1,s_2\} & \emptyset & \emptyset
  \end{array}
  \]
  This table is exactly the state-transition diagram of Figure~4.20.

  \item \textbf{Diagram of $A^{\sigma d}_\lambda$ (connected portion).}
  Starting from the actual start set $\{q_0,s_0\}$, the reachable subsets are
  \[
    R_0=\{q_0,s_0\},\quad
    R_1=\{s_1,s_2\},\quad
    R_2=\{s_3\},\quad
    R_3=\{s_2\},\quad
    R_4=\emptyset.
  \]
  Their transitions are
  \[
  \begin{array}{c|cccc}
   & \pmb{a} & \pmb{b} & \pmb{c} & \pmb{d}\\\hline
  R_0 & R_1 & R_2 & R_3 & R_3\\
  R_1 & R_4 & R_2 & R_3 & R_3\\
  R_2 & R_4 & R_1 & R_4 & R_4\\
  R_3 & R_4 & R_4 & R_4 & R_3\\
  R_4 & R_4 & R_4 & R_4 & R_4
  \end{array}
  \]
  The final states are $R_0$, $R_1$, and $R_3$.
\end{enumerate}

%% 4.32
\item \textbf{What is wrong with the given ``proof'' of Lemma~4.2?}

The flaw is in the inductive step.  The ``proof'' starts the induction at $k=1$
(base case for single letters), but the inductive step uses
$\overline{\delta^{\sigma}_\lambda}(s,x)$ for $x\in\Sigma^k$ and then applies
the ``induction hypothesis'' $\overline{\delta^{\sigma}_\lambda}(s,x)=\DBAR_\lambda(s,x)$.

The critical error is on the line:
\[
  \delta^{\sigma}_\lambda(\overline{\delta^{\sigma}_\lambda}(s,x),\pmb{a})
  = \overline{\delta^{\sigma}_\lambda}(\DBAR_\lambda(s,x),\pmb{a}).
\]
Here, $\overline{\delta^{\sigma}_\lambda}(s,x)$ is a \emph{set} of states
(since we're in an NDFA), but $\delta^{\sigma}_\lambda$ is defined on individual
states, not sets.  The expression $\delta^{\sigma}_\lambda(Q,\pmb{a})$ where
$Q$ is a set is not directly the same as $\overline{\delta^{\sigma}_\lambda}(Q,\pmb{a})$
for a set-valued argument: the extended function $\overline{\delta^{\sigma}_\lambda}$
on a set-argument must take the union over all states in the set.

More specifically, the step
\[
  \delta^{\sigma}_\lambda \bigl(\overline{\delta^{\sigma}_\lambda}(s,x),\pmb{a}\bigr)
  = \overline{\delta^{\sigma}_\lambda} \bigl(\DBAR_\lambda(s,x),\pmb{a}\bigr)
\]
conflates applying $\delta^{\sigma}_\lambda$ to a \emph{set} with applying the
set-union version of $\overline{\delta^{\sigma}_\lambda}$ to a set argument.  These
are not the same: the left side applies $\delta^{\sigma}_\lambda$ to the entire
set at once (which is not how the function is defined), while the right side
correctly takes the union over individual states.  This makes the step invalid
without justification.  The same typing problem reappears later if one writes an
expression such as
\[
\DBAR_\lambda(\DBAR_\lambda(s,x),\pmb{a}),
\]
because the first argument of $\DBAR_\lambda$ must be a single state, not a set
of states.  The correct way to proceed is to write a union over the states in
that set.

%% 4.33
\item Consider the NDFA $T$ from Example~4.7 (detecting substring $\pmb{010010}$).

\begin{enumerate}[(a)]
  \item \emph{Circuit diagram (no SOS/EOS).}  $T$ has 8 states (one start
  state that loops on all input plus the six states that track the pattern
  $\pmb{010010}$).  Using one activity bit per state, the next-state equations
  are
  \[
    s'_0=s_0,\qquad
    s'_1=s_0\wedge \pmb{0},\qquad
    s'_2=s_1\wedge \pmb{1},
  \]
  \[
    s'_3=s_2\wedge \pmb{0},\qquad
    s'_4=s_3\wedge \pmb{0},\qquad
    s'_5=s_4\wedge \pmb{1},
  \]
  \[
    s'_6=(s_5\wedge \pmb{0})\vee(s_6\wedge \pmb{0})\vee(s_6\wedge \pmb{1}).
  \]
  Initially only $s_0$ is active, and the accept line is $s_6$.

  \item \emph{DFA $T^d$ (connected portion).}  See Exercise~4.29 for the
  subset construction.  The connected DFA is exactly the 12-state machine
  listed there, with states $Q_0,\dots,Q_{11}$ and the transition table given
  in Exercise~4.29.
\end{enumerate}

%% 4.34
\item Consider the NDFA from Example~4.11, which accepts all words containing
$\pmb{10110}$, $\pmb{1010}$, or $\pmb{01101}$ as substrings.

\begin{enumerate}[(a)]
  \item \emph{Circuit diagram (no SOS/EOS).}  The NDFA has a single initial
  state $s_0$ that loops on $\Sigma^*$, three pattern-matching arms, and the
  final sink $s_{18}$ from Figure~4.12.  Since $s_1$, $s_7$, and $s_{12}$ are
  only shown for clarity in the text, they can be omitted from the actual
  circuit.  With one activity bit per remaining state, the nonzero next-state
  equations are
  \[
    s'_0=s_0,
  \]
  \[
    s'_2=s_0\wedge \pmb{1},\quad
    s'_3=s_2\wedge \pmb{0},\quad
    s'_4=s_3\wedge \pmb{1},\quad
    s'_5=s_4\wedge \pmb{1},\quad
    s'_6=s_5\wedge \pmb{0},
  \]
  \[
    s'_8=s_0\wedge \pmb{1},\quad
    s'_9=s_8\wedge \pmb{0},\quad
    s'_{10}=s_9\wedge \pmb{1},\quad
    s'_{11}=s_{10}\wedge \pmb{0},
  \]
  \[
    s'_{13}=s_0\wedge \pmb{0},\quad
    s'_{14}=s_{13}\wedge \pmb{1},\quad
    s'_{15}=s_{14}\wedge \pmb{1},\quad
    s'_{16}=s_{15}\wedge \pmb{0},\quad
    s'_{17}=s_{16}\wedge \pmb{1},
  \]
  \[
    s'_{18}=s_6\vee s_{11}\vee s_{17}\vee s_{18}.
  \]
  The accept line is $s_6\vee s_{11}\vee s_{17}\vee s_{18}$.

  \item \emph{DFA (connected portion only).}  Since all three patterns are
  started from the looping state $s_0$, the subset construction yields the
  following 23 reachable states:
  \[
  \begin{array}{lll}
  Q_0=\{s_0\}, &
  Q_1=\{s_0,s_{13}\}, &
  Q_2=\{s_0,s_2,s_8\},\\
  Q_3=\{s_0,s_{14},s_2,s_8\}, &
  Q_4=\{s_0,s_{13},s_3,s_9\}, &
  Q_5=\{s_0,s_{15},s_2,s_8\},\\
  Q_6=\{s_0,s_{10},s_{14},s_2,s_4,s_8\}, &
  Q_7=\{s_0,s_{13},s_{16},s_3,s_9\}, &
  Q_8=\{s_0,s_{11},s_{13},s_3,s_9\},\\
  Q_9=\{s_0,s_{15},s_2,s_5,s_8\}, &
  Q_{10}=\{s_0,s_{10},s_{14},s_{17},s_2,s_4,s_8\}, &
  Q_{11}=\{s_0,s_{13},s_{18}\},\\
  Q_{12}=\{s_0,s_{10},s_{14},s_{18},s_2,s_4,s_8\}, &
  Q_{13}=\{s_0,s_{13},s_{16},s_3,s_6,s_9\}, &
  Q_{14}=\{s_0,s_{11},s_{13},s_{18},s_3,s_9\},\\
  Q_{15}=\{s_0,s_{15},s_{18},s_2,s_5,s_8\}, &
  Q_{16}=\{s_0,s_{14},s_{18},s_2,s_8\}, &
  Q_{17}=\{s_0,s_{10},s_{14},s_{17},s_{18},s_2,s_4,s_8\},\\
  Q_{18}=\{s_0,s_{13},s_{16},s_{18},s_3,s_6,s_9\}, &
  Q_{19}=\{s_0,s_{18},s_2,s_8\}, &
  Q_{20}=\{s_0,s_{13},s_{18},s_3,s_9\},\\
  Q_{21}=\{s_0,s_{15},s_{18},s_2,s_8\}, &
  Q_{22}=\{s_0,s_{13},s_{16},s_{18},s_3,s_9\}. &
  \end{array}
  \]
  Their transition table is
  \[
  \begin{array}{c|cc}
   & \pmb{0} & \pmb{1}\\\hline
  Q_0 & Q_1 & Q_2\\
  Q_1 & Q_1 & Q_3\\
  Q_2 & Q_4 & Q_2\\
  Q_3 & Q_4 & Q_5\\
  Q_4 & Q_1 & Q_6\\
  Q_5 & Q_7 & Q_2\\
  Q_6 & Q_8 & Q_9\\
  Q_7 & Q_1 & Q_{10}\\
  Q_8 & Q_{11} & Q_{12}\\
  Q_9 & Q_{13} & Q_2\\
  Q_{10} & Q_{14} & Q_{15}\\
  Q_{11} & Q_{11} & Q_{16}\\
  Q_{12} & Q_{14} & Q_{15}\\
  Q_{13} & Q_{11} & Q_{17}\\
  Q_{14} & Q_{11} & Q_{12}\\
  Q_{15} & Q_{18} & Q_{19}\\
  Q_{16} & Q_{20} & Q_{21}\\
  Q_{17} & Q_{14} & Q_{15}\\
  Q_{18} & Q_{11} & Q_{17}\\
  Q_{19} & Q_{20} & Q_{19}\\
  Q_{20} & Q_{11} & Q_{12}\\
  Q_{21} & Q_{22} & Q_{19}\\
  Q_{22} & Q_{11} & Q_{17}
  \end{array}
  \]
  The final states are exactly those containing one of $s_6$, $s_{11}$,
  $s_{17}$, or $s_{18}$:
  \[
  Q_8,Q_{10},Q_{11},Q_{12},Q_{13},Q_{14},Q_{15},Q_{16},
  Q_{17},Q_{18},Q_{19},Q_{20},Q_{21},Q_{22}.
  \]
\end{enumerate}

%% 4.35
\item Consider the NDFA $B$ from Example~4.8 and its deterministic equivalent $B^d$.

\begin{enumerate}[(a)]
  \item \emph{Circuit diagram for $B$ (no SOS/EOS).}  Encode the states of $B$
  using one activity bit for $r$ and one for $s$.  From Figure~4.7,
  \[
    r'=(r\wedge \pmb{a})\vee(s\wedge \pmb{a}),\qquad
    s'=(r\wedge \pmb{a})\vee(r\wedge \pmb{b}),
  \]
  with initial state $r=1$, $s=0$.  Since $s$ is the only final state, the
  accept output is $s$.

  \item \emph{Circuit diagram for $B^d$ (no SOS/EOS).}  The states of $B^d$ are
  \[
    \emptyset,\qquad \{r\},\qquad \{s\},\qquad \{r,s\}.
  \]
  Using the binary state assignment
  \[
    \emptyset=00,\qquad \{r\}=01,\qquad \{s\}=10,\qquad \{r,s\}=11,
  \]
  with present-state bits $x_1x_0$, the transition table is
  \[
  \begin{array}{c|cc}
  x_1x_0 & \pmb{a} & \pmb{b}\\\hline
  00 & 00 & 00\\
  01 & 11 & 10\\
  10 & 01 & 00\\
  11 & 11 & 10
  \end{array}
  \]
  so the DFA circuit is given by
  \[
    x'_1=x_0,\qquad x'_0=\pmb{a}\wedge(x_1\vee x_0).
  \]
  The start state is $01$ and the accept output is $x_1$, since the final
  subset states are $\{s\}$ and $\{r,s\}$.
\end{enumerate}

%% 4.36
\item \textbf{Circuit diagrams for each NDFA in Figure~4.25 (no SOS/EOS).}

For each NDFA $A_i$ in Figure~4.25, with $\Sigma=\{\pmb{a},\pmb{b},\pmb{c}\}$:
\begin{enumerate}[1.]
  \item \textbf{(a)} With activity bit $t$ for the single state:
  \[
    t'=t\wedge(\pmb{a}\vee\pmb{c}),\qquad \text{accept}=t.
  \]

  \item \textbf{(b)} With activity bits $q_1,q_0$:
  \[
    q'_1=(q_1\wedge \pmb{b})\vee(q_0\wedge \pmb{b}),\qquad
    q'_0=(q_1\wedge \pmb{a})\vee(q_0\wedge \pmb{c}),
  \]
  and the accept output is $q_1$.

  \item \textbf{(c)} With activity bits $r_0,r_1$:
  \[
    r'_0=(r_0\wedge \pmb{a})\vee(r_1\wedge \pmb{b}),\qquad
    r'_1=(r_0\wedge \pmb{a})\vee(r_0\wedge \pmb{c}),
  \]
  and the accept output is $r_0$.

  \item \textbf{(d)} With activity bits $s_0,s_1,s_2,s_3$:
  \[
    s'_0=0,\qquad
    s'_1=(s_0\wedge \pmb{a})\vee(s_3\wedge \pmb{b}),
  \]
  \[
    s'_2=s_1\wedge \pmb{a},\qquad
    s'_3=s_1\wedge \pmb{b},
  \]
  and the accept output is $s_2$.

  \item \textbf{(e)} With activity bits $s_1,s_2$:
  \[
    s'_1=(s_1\wedge \pmb{a})\vee(s_2\wedge \pmb{a})\vee(s_2\wedge \pmb{b}),
    \qquad
    s'_2=s_1\wedge \pmb{a},
  \]
  and the accept output is $s_2$.

  \item \textbf{(f)} The $\lambda$-machine has states $s_0,s_1,s_2,s_3$ with
  ordinary transitions $s_0\xrightarrow{\pmb{a}}s_1$,
  $s_1\xrightarrow{\pmb{b}}s_2$, $s_2\xrightarrow{\pmb{c}}s_3$ and
  $\lambda$-moves $s_1\to s_0$, $s_3\to s_1$.  The standard circuit therefore
  uses
  \[
    D_0=s_1,\qquad
    D_1=(s_0\wedge \pmb{a})\vee s_3,\qquad
    D_2=s_1\wedge \pmb{b},\qquad
    D_3=s_2\wedge \pmb{c},
  \]
  with accept output $s_3$.  Adding the transitive $\lambda$-connection
  $s_3\to s_0$ as well is also consistent with the text's discussion.

  \item \textbf{(g)} With activity bits $s_0,s_1,s_2,s_3$:
  \[
    s'_0=s_1\wedge \pmb{a},\qquad
    s'_1=(s_0\wedge \pmb{a})\vee(s_2\wedge \pmb{c})\vee(s_3\wedge \pmb{b}),
  \]
  \[
    s'_2=s_1\wedge \pmb{c},\qquad
    s'_3=s_1\wedge \pmb{b},
  \]
  and the accept output is $s_0$.
\end{enumerate}

%% 4.37
\item \textbf{Circuit diagrams for each NDFA in Figure~4.25 (with SOS and EOS).}

Extend each circuit from Exercise~4.36 by forcing the start state(s) active on
\<SOS\> and gating the final-state output with \<EOS\>:
\begin{itemize}
  \item \textbf{(a)} $t'=\text{\<SOS\>}\vee(t\wedge(\pmb{a}\vee\pmb{c}))$ and
    $\text{accept}=t\wedge\text{\<EOS\>}$.
  \item \textbf{(b)} OR \<SOS\> into the input for $q_1$:
    \[
      q'_1=\text{\<SOS\>}\vee(q_1\wedge \pmb{b})\vee(q_0\wedge \pmb{b}),\qquad
      q'_0=(q_1\wedge \pmb{a})\vee(q_0\wedge \pmb{c}),
    \]
    and $\text{accept}=q_1\wedge\text{\<EOS\>}$.
  \item \textbf{(c)} OR \<SOS\> into the input for $r_0$ and gate the output:
    \[
      r'_0=\text{\<SOS\>}\vee(r_0\wedge \pmb{a})\vee(r_1\wedge \pmb{b}),\qquad
      r'_1=(r_0\wedge \pmb{a})\vee(r_0\wedge \pmb{c}),
    \]
    with $\text{accept}=r_0\wedge\text{\<EOS\>}$.
  \item \textbf{(d)} OR \<SOS\> into $s'_0$ and use
    $\text{accept}=s_2\wedge\text{\<EOS\>}$.
  \item \textbf{(e)} OR \<SOS\> into the corrected first equation:
    \[
      s'_1=\text{\<SOS\>}\vee(s_1\wedge \pmb{a})
        \vee(s_2\wedge \pmb{a})\vee(s_2\wedge \pmb{b}),
      \qquad
      s'_2=s_1\wedge \pmb{a},
    \]
    and use $\text{accept}=s_2\wedge\text{\<EOS\>}$.
  \item \textbf{(f)} OR \<SOS\> into $D_0$ and use
    $\text{accept}=s_3\wedge\text{\<EOS\>}$.
  \item \textbf{(g)} OR \<SOS\> into $s'_0$ and use
    $\text{accept}=s_0\wedge\text{\<EOS\>}$.
\end{itemize}

%% 4.38
\item \textbf{Adapting Definition~3.10 for NDFAs and proving that an algorithm
exists to find the connected portion.}

\begin{enumerate}[(a)]
  \item \textbf{Adapted definition.}  The connected portion of an NDFA
  $A=\langle\Sigma,S,S_0,\delta,F\rangle$ is defined as follows.  Let
  \[
    \text{Reach}_0 = S_0,\qquad
    \text{Reach}_{k+1} = \text{Reach}_k
      \cup\bigcup_{\pmb{a}\in\Sigma}\bigcup_{s\in\text{Reach}_k}\delta(s,\pmb{a}).
  \]
  The connected portion is $\text{Reach}=\bigcup_{k\ge 0}\text{Reach}_k$, together
  with all transitions among states in $\text{Reach}$ and the restriction of
  $F$ to $\text{Reach}$.

  \item \textbf{Proof that the algorithm terminates.}

  \begin{proof}
  The sets $\text{Reach}_k$ form a non-decreasing chain
  $\text{Reach}_0\subseteq\text{Reach}_1\subseteq\cdots\subseteq S$.
  Since $S$ is finite, the chain must stabilize: there exists $k^*$ such that
  $\text{Reach}_{k^*}=\text{Reach}_{k^*+1}$.  At that point, no new states are
  added.  The algorithm terminates in at most $|S|$ iterations (since at least
  one new state must be added in each non-stable step, and there are only $|S|$
  states).

  To compute this algorithmically:
  \begin{enumerate}[1.]
    \item Initialize $R:=S_0$, $\text{New}:=S_0$.
    \item While $\text{New}\neq\emptyset$:
      \begin{enumerate}[(a)]
        \item $\text{Next}:=\bigcup_{\pmb{a}\in\Sigma}\bigcup_{s\in\text{New}}
              \delta(s,\pmb{a})\setminus R$.
        \item $R:=R\cup\text{Next}$; $\text{New}:=\text{Next}$.
      \end{enumerate}
    \item Output $R$ as the connected portion.
  \end{enumerate}
  Each iteration of the while loop adds at least one state to $R$ (otherwise
  $\text{New}=\emptyset$ and we stop).  Thus the loop executes at most $|S|$
  times.  The running time depends on how the transition sets are represented:
  with adjacency lists it is linear in the number of reachable transition
  triples $(s,\pmb{a},t)$ with $t\in\delta(s,\pmb{a})$, while a dense
  subset-valued representation gives a bound of $O(|S|^2\cdot|\Sigma|)$.
  In either representation the procedure is finite and effective. \qed
  \end{proof}
\end{enumerate}

\end{enumerate}
